\documentclass[a4paper,12pt]{article}
\usepackage[T2A]{fontenc}
\usepackage[utf8]{inputenc}
\usepackage[russian]{babel}
\usepackage{amsmath}
\usepackage{amssymb}
\usepackage[margin=2cm]{geometry}

\begin{document}

\begin{center}
    {\Large Построение ДКА по регулярному выражению на основе алгоритма Глушкова}
\end{center}

\section*{Постановка}

На вход подаётся регулярное выражение. Надо:

\begin{enumerate}
    \item разобрать его;
    \item построить по нему автомат на основе алгоритма Глушкова;
    \item получить ДКА;
    \item сгенерировать latex код его диаграммы.
\end{enumerate}

Язык регулярных выражений:

\[
R ::= 0 \mid 1 \mid a \mid (R \mid R) \mid RR \mid R^*
\]

$a$ --- символ алфавита, $0$ обозначает пустой язык, $1$ --- пустое слово.

\section*{Алгоритм}

Регулярное выражение линеаризуется -- все вхождения символа получают номер.
Например $(a \mid b)^*abb$ ->$(a_1 \mid b_2)^* a_3 b_4 b_5.$
Дальше вычисляем:

\begin{enumerate}
    \item $\text{nullable}(R)$ --- допускает ли выражение пустое слово;
    \item $\text{firstpos}(R)$ --- позиции, с которых может начинаться слово;
    \item $\text{lastpos}(R)$ --- позиции, которыми слово может заканчиваться;
    \item $\text{followpos}(i)$ --- позиции, которые могут идти сразу после позиции $i$.
\end{enumerate}

\section*{Рекурсивные правила}

Для объединения:

\[
\text{firstpos}(R_1 \mid R_2) = \text{firstpos}(R_1) \cup \text{firstpos}(R_2)
\]

\[
\text{lastpos}(R_1 \mid R_2) = \text{lastpos}(R_1) \cup \text{lastpos}(R_2)
\]

Для конкатенации:

\[
\text{firstpos}(R_1R_2) =
\begin{cases}
\text{firstpos}(R_1) \cup \text{firstpos}(R_2), & \text{если } \text{nullable}(R_1) \\
\text{firstpos}(R_1), & \text{иначе}
\end{cases}
\]

\[
\text{lastpos}(R_1R_2) =
\begin{cases}
\text{lastpos}(R_1) \cup \text{lastpos}(R_2), & \text{если } \text{nullable}(R_2) \\
\text{lastpos}(R_2), & \text{иначе}
\end{cases}
\]

И для каждой позиции из $\text{lastpos}(R_1)$ в множество $\text{followpos}$ добавляются все позиции из $\text{firstpos}(R_2)$.

Для звезды Клини $R^*$:

\[
\text{nullable}(R^*) = \text{True}
\]

\[
\text{firstpos}(R^*) = \text{firstpos}(R)
\]

\[
\text{lastpos}(R^*) = \text{lastpos}(R)
\]

И для каждой позиции из $\text{lastpos}(R)$ в $\text{followpos}$ добавляются все позиции из $\text{firstpos}(R)$.

\section*{Построение ДКА}

Сначала по вычисленным множествам задаётся автомат Глушкова:

\begin{enumerate}
    \item начальное состояние соответствует началу слова;
    \item состояния, помеченные последними позициями, принимающие;
    \item переходы определяются через $\text{firstpos}$ и $\text{followpos}$.
\end{enumerate}

После этого строится ДКА методом подмножеств.
Состояние ДКА -- множество позиций исходного позиционного автомата.

Начальное состояние ДКА: $\varnothing$

Оно является принимающим $\longleftrightarrow$ исходное регулярное выражение допускает пустое слово.

\section*{Результат работы}

Программа печатает:

\begin{enumerate}
    \item список размеченных позиций
    \item множества $FirstPos$, $LastPos$, $FollowPos$
    \item список состояний ДКА
    \item список переходов ДКА.
\end{enumerate}

И программа создаёт .tex - файл с TikZ диаграммой автомата.

\end{document}